\section{Программный интерфейс C++}
\label{sec:api}

\subsection{Публичные функции управления}
Основной заголовок подключения --- \code{back/simulation.h}.
\begin{lstlisting}[language=C++]
void configure_simulation(const SimulationParameters& parameters);
const SimulationParameters& simulation_parameters();
const SimulationResult& simulation_result();
void init_simulation();
void manage_warehouse(int current_day);
void generate_daily_events(int current_day);
void process_daily_order(int current_day);
void cleanup_random_customers();
void print_final_statistics();
\end{lstlisting}

\begin{tabularx}{\textwidth}{@{}L{0.39\textwidth}Y@{}}
\toprule
Функция & Контракт \\
\midrule
\code{configure_simulation} & Проверяет и сохраняет параметры. Допустима
до инициализации или после завершения дня; внутри дня запрещена.
Для запуска с новыми параметрами затем нужна инициализация. \\
\code{simulation_parameters} & Возвращает константную ссылку на текущие параметры. \\
\code{init_simulation} & Создаёт начальное состояние и сбрасывает результаты,
очереди и статистику эксперимента. \\
\code{manage_warehouse} & Начинает следующий день: склад, потери, поставки. \\
\code{generate_daily_events} & Создаёт заказы текущего начатого дня. \\
\code{process_daily_order} & Доставляет вчерашнее, готовит завтрашнее,
формирует пополнение. \\
\code{cleanup_random_customers} & Завершает день, обновляет покупателей
и добавляет дневной результат. \\
\code{simulation_result} & Возвращает константную ссылку на результат;
актуальный снимок готов после завершения дня. \\
\code{print_final_statistics} & Выводит накопленную итоговую статистику
в \code{std::cout}. \\
\bottomrule
\end{tabularx}

Возвращаемые ссылки относятся к единственному экземпляру движка.
Чтобы сохранить результаты независимо от следующих дней или нового запуска,
скопируйте \code{SimulationResult}. Все вызовы для одного эксперимента должны
выполняться последовательно. Публичная функция \code{init_simulation()}
сбрасывает модель; не используйте её как шаг продолжения текущего дня.

\subsection{Пример подключения модели}
Минимальный пример потребителя, выполняющего все дни без Qt-точки входа:
\begin{lstlisting}[language=C++]
#include "back/simulation.h"
#include <exception>
#include <iostream>

int main() {
    try {
        SimulationParameters p;
        p.days = 10;
        p.couriers = 5;
        p.medicine_count = 15;
        p.markup_percent = 25.0;
        p.seed = 42;

        configure_simulation(p);
        init_simulation();
        for (int day = 1; day <= p.days; ++day) {
            manage_warehouse(day);
            generate_daily_events(day);
            process_daily_order(day);
            cleanup_random_customers();
            const auto& current = simulation_result().days.back();
            std::cout << current.day << ' '
                      << current.delivered_orders << '\n';
        }
        SimulationResult saved = simulation_result();
        std::cout << saved.profit << '\n';
        return 0;
    } catch (const std::exception& error) {
        std::cerr << error.what() << '\n';
        return 1;
    }
}
\end{lstlisting}
В приложении с графическим интерфейсом тот же порядок можно выполнять
по пользовательской команде или таймеру. Обновляйте представление после
\code{cleanup_random_customers()}. При отдельной сборке потребителя подключите
реализации склада, заказа, статистики, генератора, контекста, поставок и движка;
\code{main.cpp} консольного приложения в такую сборку не включается.

\subsection{Параметры и сущности}
\code{SimulationParameters} содержит поля \code{days}, \code{couriers},
\code{medicine_count}, \code{markup_percent}, \code{card_discount},
\code{large_purchase_discount}, \code{regular_discount},
\code{maximum_discount}, \code{initial_count}, \code{regular_customers},
\code{seed} и \code{initial_stock_file}. Значения по умолчанию и ограничения
приведены в разделе~\ref{sec:running}; \code{validate()} проверяет их
без запуска эксперимента. Соответствие имён CLI и полей:
\code{--bulk-discount} задаёт \code{large_purchase_discount},
\code{--stock} задаёт \code{initial_stock_file}.

{\small
\begin{longtable}{@{}L{0.21\textwidth}L{0.72\textwidth}@{}}
\toprule
Сущность & Поля и смысл \\
\midrule
\endhead
\code{Medicine} & \code{dosage}, \code{type_med_id}, \code{group_med_id},
\code{wholesale_price}, \code{shelf_life}, \code{min_stock}:
дозировка, индексы формы и группы, закупочная цена, срок годности и порог пополнения. \\
\code{Customer} & \code{fio}, \code{phone}, \code{address};
\code{has_card}, \code{is_regular}, \code{card_number};
\code{regular_purchases}, \code{periodicity}, \code{next_purchase_day}.
В \code{regular_purchases} хранятся пары «индекс лекарства, количество». \\
\code{StoreBatch} & \code{mas_med_id}, \code{count}, \code{expiration_day},
\code{batch_price}: лекарство, текущий остаток, последний день годности и
текущая продажная база за упаковку. Закупочная цена хранится отдельно. \\
\code{Warehouse} & \code{inventory}: вектор партий. Количество одного лекарства
получается суммированием нескольких партий. \\
\code{Order} & \code{customer_idx}, \code{items}, \code{final_price}:
индекс покупателя, исходный запрос и рассчитанная цена фактической выдачи. \\
\code{Request} & \code{mas_med_id}, \code{count}, \code{delivery_timer}:
лекарство, количество и оставшееся время поставки. \\
\code{Statistics} & \code{income}, \code{purchase_expenses},
\code{write_off_losses}: накопленные выручка, закупки и потери. \\
\bottomrule
\end{longtable}
}

Публичные числовые поля исходных классов не имеют значений по умолчанию.
При самостоятельном создании используйте инициализацию через фигурные скобки,
например \code{Customer customer{};} и \code{Statistics stats{};}.
Индекс лекарства должен соответствовать текущему \code{mas_med};
индекс клиента --- текущему \code{customers}.

\subsection{Методы предметных классов}
\begin{description}
  \item[\code{StoreBatch::price()}] Возвращает \code{count * batch_price},
  то есть стоимость остатка по текущей базе. Проверяет неотрицательность
  количества и цены, конечность результата.
  \item[\code{StoreBatch::take(int cnt)}] Уменьшает остаток на
  \code{min(cnt, count)} и возвращает реально взятое. Отрицательные
  количества отклоняются.
  \item[\code{Warehouse::add_batch(const StoreBatch&)}] Проверяет идентификатор,
  количество, срок и цену; добавляет партию и её закупочные сведения.
  Метод не начисляет расходы: этим занимается вызывающая логика поставки
  или инициализации.
  \item[\code{Warehouse::write_off(int day)}] Удаляет истёкшие и пустые партии,
  применяет однократную уценку. Потери сохраняются во внутреннем контексте;
  сбор и перенос потерь в итоговую статистику выполняет движок.
  \item[\code{Warehouse::need_restock()}] Проверяет, существует ли хотя бы
  одно лекарство с суммой действительных остатков ниже \code{min_stock}.
  Сам по себе метод не создаёт заявок.
  \item[\code{Order::process_order(Warehouse&, Statistics&, double)}]
  Проверяет запрос, рассчитывает выдачу и цену, снимает упаковки, обновляет
  переданную статистику и внутренний учёт. Возвращает \code{true}, если
  выдана хотя бы одна упаковка; полнота запроса этим значением не гарантируется.
  \item[\code{Statistics::update(double, double, double)}] Прибавляет выручку,
  закупки и потери; проверяет конечность, неотрицательность исходных и новых сумм.
\end{description}

Низкоуровневый \code{Order::process_order} сразу добавляет цену в переданную
статистику. Движок вызывает его с временной \code{Statistics reserved_sale{}},
чтобы итоговую выручку начислить только при доставке на следующий день.
При прямом вызове дата и допустимый срок доставки берутся из
\code{ModelContext}; стандартный дневной контракт устанавливает их сам.

Внутреннее \code{last_fulfillment} содержит фактическую выдачу только
последней обработки и заменяется следующей. Для внешнего потребителя
устойчивым источником истории служат \code{CompletedPurchase}.

\subsection{Выполненная покупка}
\code{CompletedPurchase} содержит:
\begin{itemize}
  \item \code{order_id}: последовательный идентификатор с 1, сбрасывается
        при новом эксперименте;
  \item \code{ordered_day}, \code{prepared_day}, \code{delivered_day}:
        день получения, выделения и доставки;
  \item \code{courier_index}: индекс назначенного курьера;
  \item \code{customer}: копию всех полей \code{Customer};
  \item \code{requested}, \code{supplied}: векторы пар «лекарство, упаковки»;
  \item \code{price}: окончательную цену;
        \code{acquisition_cost}: закупочную стоимость выданных упаковок.
\end{itemize}

\subsection{Результат дня и всего эксперимента}
\code{DailySimulationResult} хранит \code{day}, \code{received_orders},
\code{planned_orders}, \code{empty_orders}, \code{prepared_orders},
\code{delivered_orders}, \code{late_deliveries}, \code{arrived_supplies},
\code{created_supplies}, \code{pending_orders},
\code{courier_limits_satisfied}, \code{courier_load},
\code{available_packages} и накопленную \code{totals}.
Счётчик \code{planned_orders} входит в \code{received_orders}.
Размеры массивов нагрузки и остатков равны $M$ и $K$ соответственно.

\code{SimulationResult} содержит \code{totals}, \code{profit},
\code{available_stock_cost}, \code{reserved_stock_cost},
\code{delivered_stock_cost}, \code{initial_packages},
\code{supplied_packages}, \code{written_off_packages},
\code{pending_deliveries}, \code{pending_orders}, \code{pending_supplies},
\code{days} и \code{completed_orders}. Различайте:
\begin{itemize}
  \item \code{pending_orders}: запросы в очереди, ещё без выделенного товара;
  \item \code{pending_deliveries}: уже подготовленные покупки на следующий день;
  \item \code{pending_supplies}: ожидающие заявки поставщикам.
\end{itemize}
Полный смысл сериализованных полей приведён в разделе~\ref{sec:formats}.

\subsection{Исключения и восстановление}
\code{std::invalid_argument} сигнализирует о недопустимых параметрах и полях,
\code{std::logic_error} --- о нарушении последовательности или внутренних
инвариантов, \code{std::overflow_error} --- о числовом переполнении.
Файловые ошибки обычно представлены \code{std::runtime_error}; разбор
стандартными \code{stoll}, \code{stod} и \code{stoull} также может выдавать
\code{std::out_of_range}.

Интерфейс не гарантирует откат целого дня при исключении. Если шаг завершился
ошибкой, восстановите корректные параметры и начните новый эксперимент
через инициализацию; продолжать частично выполненную фазу нельзя считать
поддерживаемым способом восстановления.
